\ticket{7}{Сетевая модель представления знаний. Семантические сети и их свойства. Фреймы, их виды и структура, межфреймовые связи. Представление значений по умолчанию, понятие немонотонного вывода. Графы знаний: концепция, примеры}

\begin{examplan}
    \item Описать семантические сети и их свойства.
    \item Раскрыть фреймы, их структуру, виды и межфреймовые связи.
    \item Объяснить значения по умолчанию и немонотонный вывод.
    \item Дать понятие графа знаний и привести примеры.
\end{examplan}

\subsection*{Короткая версия для письменного ответа}

\textbf{Сетевая модель представления знаний} описывает знания в виде графа: вершины соответствуют объектам, понятиям или событиям, а ребра --- отношениям между ними. К сетевым моделям относят семантические сети, фреймовые системы и современные графы знаний.

\subsection*{Семантические сети}

\textbf{Семантическая сеть} --- ориентированный граф, где узлы обозначают понятия или объекты, а дуги с метками обозначают семантические отношения.

Типичные отношения:
\begin{itemize}
    \item \textbf{IS-A}: экземпляр класса или подкласс класса;
    \item \textbf{PART-OF}: часть-целое;
    \item \textbf{HAS-A}: обладание или наличие свойства;
    \item причинные, временные, пространственные и агентивные отношения.
\end{itemize}

Главные свойства семантических сетей:
\begin{itemize}
    \item наглядность;
    \item ассоциативность;
    \item наследование свойств по отношениям типа IS-A;
    \item гибкость расширения;
    \item вывод через поиск путей и наследование.
\end{itemize}

Ограничения: неоднозначность меток дуг, слабая формальная семантика, трудности с кванторами, отрицанием и сложными логическими выражениями.

\subsection*{Фреймы}

\textbf{Фрейм} --- структура для представления типичного объекта, понятия или ситуации через набор слотов и значений.

Фрейм обычно включает:
\begin{itemize}
    \item имя фрейма;
    \item слоты, то есть атрибуты;
    \item значения слотов;
    \item фасеты слотов: \texttt{VALUE}, \texttt{DEFAULT}, \texttt{IF-NEEDED}, \texttt{IF-ADDED}, тип, диапазон, кардинальность.
\end{itemize}

Значением слота может быть простое значение, ссылка на другой фрейм или процедура-демон, вызываемая при обращении к слоту.

\subsection*{Виды и связи фреймов}

\begin{itemize}
    \item \textbf{Фреймы-прототипы} описывают класс объектов или типовую ситуацию.
    \item \textbf{Фреймы-экземпляры} описывают конкретный объект или случай.
\end{itemize}

Межфреймовые связи:
\begin{itemize}
    \item \textbf{IS-A} или \textbf{AKO}: связь с более общим фреймом, наследование слотов;
    \item ссылки через значения слотов;
    \item связи потомков и экземпляров.
\end{itemize}

\subsection*{Значения по умолчанию и немонотонный вывод}

Фреймы позволяют задавать значения по умолчанию. Если у экземпляра значение не указано, оно наследуется от прототипа.

\textbf{Немонотонный вывод} --- вывод, который может быть отменен при добавлении новой информации. Пример: по умолчанию птицы летают; если известно только, что Твити --- птица, делается вывод, что Твити летает. Если затем выясняется, что Твити --- пингвин, вывод отменяется.

Такой вывод нужен для моделирования здравого смысла и рассуждений при неполной информации.

\subsection*{Графы знаний}

\textbf{Граф знаний} --- структурированное представление фактов о сущностях, их атрибутах и отношениях в виде графа. Часто факты задаются тройками: субъект, предикат, объект.

Пример тройки:
\begin{verbatim}
(Leonardo_da_Vinci, painted, Mona_Lisa)
(Mona_Lisa, located_in, Louvre)
\end{verbatim}

Характерные черты графов знаний:
\begin{itemize}
    \item сущности имеют идентификаторы;
    \item отношения типизированы;
    \item могут использоваться RDF, RDFS, OWL;
    \item поддерживаются интеграция данных, семантический поиск, рекомендации и вопросно-ответные системы.
\end{itemize}

Графы знаний можно рассматривать как современное развитие идей семантических сетей с большим масштабом и более формальными стандартами.

\begin{important}
    \item Семантическая сеть представляет знания как граф понятий и отношений.
    \item Фрейм группирует знания об объекте или ситуации в слотах.
    \item Наследование значений по умолчанию является примером немонотонного вывода.
    \item Граф знаний хранит факты о сущностях в виде связанного графа, часто через тройки.
    \item Главное различие: семантические сети исторически более неформальны, графы знаний масштабнее и чаще стандартизованы.
\end{important}
